\documentclass[11pt,a4paper]{article}
\usepackage[T1]{fontenc}
\usepackage[utf8]{inputenc}
\usepackage{mathpazo}
\usepackage{courier}
\usepackage{amsmath,amssymb,bm}
\usepackage{booktabs,array,tabularx}
\usepackage[margin=2.6cm]{geometry}
\usepackage{microtype}
\usepackage{xcolor}
\usepackage{listings}
\usepackage{upquote}
\usepackage[authoryear,round]{natbib}
\usepackage[colorlinks=true,linkcolor=blue!50!black,citecolor=blue!50!black,urlcolor=blue!50!black]{hyperref}

\lstset{basicstyle=\ttfamily\footnotesize,columns=fullflexible,keepspaces=true,
  frame=single,framerule=0.3pt,rulecolor=\color{black!30},
  backgroundcolor=\color{black!3},xleftmargin=0.6em,xrightmargin=0.6em,
  aboveskip=0.8em,belowskip=0.8em,breaklines=true}
\newcommand{\cmd}[1]{\texttt{#1}}
\newcommand{\yv}{\mathbf y}
\newcommand{\xv}{\mathbf x}
\newcommand{\wv}{\mathbf w}
\newcommand{\zv}{\mathbf z}
\newcommand{\E}{\mathrm E}
\setlength{\parskip}{0.35em}
\setlength{\emergencystretch}{3em}

\title{\textbf{hmdbivpois}: hidden Markov bivariate Poisson regression\\ with informative dropout in Stata\\[0.4em]
\large Technical guide, version 1.0.0}
\author{Paolo Li Donni\\ University of Palermo\\ \href{mailto:paolo.lidonni@unipa.it}{paolo.lidonni@unipa.it}}
\date{November 2026}

\begin{document}
\maketitle

\begin{abstract}
\noindent This guide documents the model, the estimation method and the postestimation tools of
\cmd{hmdbivpois}, a Stata command for the hidden Markov dropout bivariate Poisson model of
\citet{carrieri2026}. The command models two counts observed on the same individuals over time when
individuals may leave the panel, for example by death, and the dropout is informative. A latent state
that follows a first-order Markov chain and a time-constant latent class shift the means of the two
counts, which have a bivariate Poisson distribution of the form of \citet{lakshminarayana1999} and
\citet{famoye2010}, and the probability of leaving the panel after each period. The dropout is missing
not at random through the shared latent variables. Estimation is by the EM algorithm with scaled
forward--backward recursions on the joint chain of classes and states, followed by direct
maximization when needed. The guide also reports Monte Carlo evidence. The command, its help file and
an example dataset are available at \url{https://github.com/plidonni/hmdbivpois}.
\end{abstract}

\tableofcontents
\bigskip

%=====================================================================
\section{Introduction}
%=====================================================================

Panels of health care use often end with death. In the application of \citet{carrieri2026}, about
30\,000 Danish patients with type-2 diabetes are followed quarterly for five years after an emergency
admission: outpatient and emergency secondary care are two related counts, and many patients die
during the follow-up. The patients who leave are those whose use of care was high and rising, so the
dropout is informative: a model estimated on the survivors, as if the dropout were at random, gives
biased estimates of the dynamics of care.

\citet{laudicella2022} model two counts with a latent state that follows a Markov chain (time-varying
unobserved heterogeneity), lagged counts (state dependence) and a bivariate Poisson distribution
within the state (contemporaneous dependence); \cmd{hmbivpois} implements that model.
\citet{carrieri2026} add a time-constant latent class, a persistent frailty, and an equation for the
dropout that shares the latent variables with the counts, in the spirit of the shared-parameter models
of \citet{bartolucci2015}. \cmd{hmdbivpois} implements this model. Without the dropout equation and with
one class it reduces to the model of \cmd{hmbivpois}, and gives the same estimates.

Section~\ref{sec:model} presents the model, Section~\ref{sec:estimation} the likelihood and its
estimation, Section~\ref{sec:post} inference and postestimation, Section~\ref{sec:command} the syntax
with examples, and Section~\ref{sec:mc} Monte Carlo evidence.

%=====================================================================
\section{The model}\label{sec:model}
%=====================================================================

\subsection{Data and notation}

Individuals are indexed by $i=1,\dots,n$ and periods by $t=0,1,\dots,T_i$. In each period we observe two
counts $\yv_{it}=(y_{1it},y_{2it})$, regressors $\xv_{1it}$, $\xv_{2it}$, $\xv_{dit}$ for the three
equations and the dropout indicator $D_{it}$, equal to one if the individual leaves the panel after
period $t$. $D_{it}$ can be one only in the last observed period $T_i$. The first observed period,
$t=0$, is the \emph{presample period}: it supplies the lagged counts of period~1 and the presample
counts $\yv_{i0}$ that enter the prior probabilities of the latent variables, and it is not an
estimation period. The periods of an individual must be consecutive.

An individual is \emph{at risk} of dropout in all the estimation periods except the last one when
$D_{iT_i}=0$: there the follow-up ends, or the individual is censored, and the dropout is not
observed. Let $r_{it}$ be the indicator of the periods at risk. The time of the dropout is thus
interval censored at the length of a period, as in \citet[eq.~6]{carrieri2026}.

Two latent variables describe the unobserved heterogeneity:
\begin{itemize}
\item a \emph{class} $C_i\in\{1,\dots,k_1\}$, constant over time (option \cmd{classes()}), a persistent
frailty;
\item a \emph{state} $V_{it}\in\{1,\dots,k_2\}$ that follows a first-order Markov chain (option
\cmd{k()}), a health status that can change over time.
\end{itemize}
$C_i$ and $V_{it}$ are independent a priori.

\subsection{Measurement model of the counts}

Given $C_i=u$ and $V_{it}=v$, the conditional means of the two counts are
\begin{equation}\label{eq:mean}
\log\mu_{jit}(u,v)=\alpha_{jv}+\delta_{ju}+\xv_{jit}\bm\beta_j+\xv^S_{jit}\bm\beta_{jv}
+\gamma_{j1}\,g(y_{1i,t-1})+\gamma_{j2}\,g(y_{2i,t-1}),\qquad j=1,2,
\end{equation}
with $\delta_{j1}=0$. The state intercepts $\alpha_{jv}$ take the place of a constant; the class shifts
$\delta_{ju}$ add to them, so that classes and states have additive effects on the log means, as in
\citet{carrieri2026}. $\bm\beta_j$ are the slopes of the regressors, $g(\cdot)$ a transformation of the
lagged counts (Section~\ref{sec:lags}). The regressors $\xv^S_{jit}$ listed in \cmd{statevary()}, and the
lags listed in \cmd{lagvary()}, have coefficients that depend on the state (not on the class); they
are reported in the equations \textit{depvar}\cmd{\_s1}, \dots, \textit{depvar}\cmd{\_s}$k_2$.

Given the latent variables and the lags, the two counts have the bivariate Poisson distribution of
\citet{lakshminarayana1999} and \citet{famoye2010},
\begin{equation}\label{eq:famoye}
p(y_1,y_2)=\frac{\mu_1^{y_1}e^{-\mu_1}}{y_1!}\,\frac{\mu_2^{y_2}e^{-\mu_2}}{y_2!}
\Big[1+\lambda\,\phi_1(y_1)\,\phi_2(y_2)\Big],
\qquad \phi_j(y)=e^{-y}-e^{-q\mu_j},\quad q=1-e^{-1}.
\end{equation}
The margins are exactly Poisson, because $\E(e^{-Y})=e^{-q\mu}$ for $Y\sim\mathrm{Po}(\mu)$; the
covariance is $\lambda q^2\mu_1\mu_2e^{-q(\mu_1+\mu_2)}$, with the sign of $\lambda$. The probability is
nonnegative only if $\lambda$ lies in an interval that depends on the means,
\begin{equation}\label{eq:range}
-\frac{1}{\max\{d_1d_2,\,(1-d_1)(1-d_2)\}}\;\le\;\lambda\;\le\;
\frac{1}{\max\{(1-d_1)d_2,\,d_1(1-d_2)\}},\qquad d_j=e^{-q\mu_j};
\end{equation}
the command keeps $\lambda$ in the intersection of these intervals over the observations and the latent
types, and stores it in \cmd{e(lambda\_range)}.

\subsection{Dropout}

The probability of leaving the panel after period $t$ is a logit with the same latent variables,
\begin{equation}\label{eq:drop}
\operatorname{logit}\Pr(D_{it}=1\mid u,v)=\alpha_{dv}+\delta_{du}+\xv_{dit}\bm\beta_d+\xv^S_{dit}\bm\beta_{dv}
+\gamma_{d1}\,g(y_{1i,t-1})+\gamma_{d2}\,g(y_{2i,t-1}),
\end{equation}
with $\delta_{d1}=0$. With \cmd{dropcounts(current)} the counts of the period, $g(y_{1it})$ and
$g(y_{2it})$, which are observed before the dropout, replace the lagged ones; with
\cmd{dropcounts(both)} both enter. In the application, the counts of the quarter before death capture
the use of care at the end of life.

\paragraph{Missing not at random.} Given the latent variables, the regressors and the observed counts,
the dropout does not depend on the counts that would have been observed after it. The dependence of the
dropout on the unobserved future counts passes entirely through $C_i$ and the sequence of $V_{it}$: the
model is a shared-parameter model, and the dropout is informative whenever the latent variables shift
both the counts and the dropout.

\paragraph{Classes in selected equations.} With \cmd{classeq(dropout)} the class shifts of the counts are
zero, $\delta_{1u}=\delta_{2u}=0$: the class is a frailty that affects only the risk of dropout. With
\cmd{classeq(counts)}, $\delta_{du}=0$. Both are nested in the default model with the restricted
parameters inside the parameter space, so a likelihood-ratio test against the default has the usual
$\chi^2$ distribution, with $2(k_1-1)$ and $k_1-1$ degrees of freedom.

\subsection{Transformations of the lagged counts}\label{sec:lags}

The option \cmd{lagtransform()} sets $g(\cdot)$ for the lags in the three equations, for the counts of
the period in the dropout equation and for the presample counts in the prior probabilities:

\begin{center}
\begin{tabular}{ll}
\toprule
Option & $g(y)$\\
\midrule
\cmd{linear} (default) & $y$\\
\cmd{log1p} & $\log(1+y)$\\
\cmd{logc(\#)} & $\log\max(y,\#)$, $0<\#\le1$\\
\cmd{logzero} & $\log y$ if $y>0$, $0$ otherwise, plus the indicator $\mathbf 1(y=0)$\\
\cmd{indicator} & $\mathbf 1(y>0)$\\
\cmd{topcode(\#)} & $\min(y,\#)$\\
\bottomrule
\end{tabular}
\end{center}

With linear lags and positive own-lag coefficients the conditional mean grows exponentially in the past
count, and the dynamics can be explosive; the log-type transformations make the model a log-linear
autoregression in the counts \citep{blundell2002}. \cmd{estat stability} diagnoses both cases.

\subsection{Prior probabilities and transitions}

\paragraph{Classes.} The prior probability of class $u$ is a multinomial logit on the covariates
$\zv_{i}$ of \cmd{classcov()}, the transformed presample counts and a constant, with class~1 as the base
(equations \cmd{class2}, \dots).

\paragraph{Initial states.} The probability of state $v$ in the first estimation period is a multinomial
logit on the covariates of \cmd{initcov()}, the transformed presample counts and a constant, with
state~1 as the base (equations \cmd{init2}, \dots), following \citet{heckman1981} and
\citet{wooldridge2005} for the initial conditions.

\paragraph{Transitions.} Each row of the transition matrix of the states is a multinomial logit with
staying in the state as the base,
\begin{equation}\label{eq:trans}
\pi_{b\mid a,it}=\Pr(V_{it}=b\mid V_{i,t-1}=a)=
\frac{\exp(\tau_{ab}+\wv_{it}\bm\psi_{ab})}{1+\sum_{h\ne a}\exp(\tau_{ah}+\wv_{it}\bm\psi_{ah})},
\qquad b\ne a,
\end{equation}
reported in the equations \cmd{tr}$a$\cmd{\_}$b$: homogeneous by default (no $\wv_{it}$); with covariates
measured at the period of arrival (\cmd{xtrans()}); with a free matrix for each period of arrival
(\cmd{tvtrans}); or without transitions (\cmd{notransition}), in which case the states, like the
classes, are constant. The classes never change, and the transition probabilities do not depend on
the class.

\subsection{Number of parameters}

With $p_j$ common slopes and $s_j$ state-specific slopes in equation $j\in\{1,2,d\}$, $m_v$ covariates
of the initial states, $m_u$ of the classes, $L$ features of the presample counts and $r$ covariates of
the transitions, the model has
\[
\begin{split}
&\sum_{j\in\{1,2,d\}}\big(p_j+k_2+e_j(k_1-1)+k_2s_j\big)+1\\
&\quad+(k_2-1)(m_v+L+1)+(k_1-1)(m_u+L+1)+k_2(k_2-1)(1+r)
\end{split}
\]
parameters, where $e_j=1$ if the classes shift equation $j$ (\cmd{classeq()}) and the 1 is $\lambda$.

%=====================================================================
\section{Likelihood and estimation}\label{sec:estimation}
%=====================================================================

\subsection{The joint chain of classes and states}

Let $f_{it}(u,v)=p(\yv_{it}\mid u,v)\,\big[p_{dit}(u,v)^{D_{it}}\{1-p_{dit}(u,v)\}^{1-D_{it}}\big]^{r_{it}}$
be the contribution of period $t$ given the latent variables, with $p_{dit}$ from \eqref{eq:drop}. The
likelihood of individual $i$ is
\begin{equation}
L_i=\sum_{u=1}^{k_1}\Pr(C_i=u)\sum_{v_1,\dots,v_{T_i}}\Pr(V_{i1}=v_1)\,f_{i1}(u,v_1)
\prod_{t=2}^{T_i}\pi_{v_t\mid v_{t-1},it}\,f_{it}(u,v_t),
\qquad\ell(\bm\vartheta)=\sum_{i=1}^n\log L_i.
\end{equation}
The command treats the pair $(C_i,V_{it})$ as a single Markov chain on $K=k_1k_2$ \emph{latent types},
$s=(u-1)k_2+v$, with initial probabilities $\Pr(C_i=u)\Pr(V_{i1}=v)$ and the block-diagonal transition
matrix $\mathbf I_{k_1}\otimes\bm\Pi$: the class never changes. This is equivalent to the recursions by
class of \citet{carrieri2026} and lets the likelihood be computed by the forward recursion of hidden
Markov models \citep{baum1970,zucchini2016} on the joint chain. The recursion is scaled at each period to
avoid underflow, and vectorized over individuals, who are sorted by decreasing $T_i$, so that the
unbalanced panels created by the dropout need no special treatment. The cost grows with $k_1^2k_2^2$,
which is negligible for the small numbers of classes used in practice.

The intercepts of the types are linear in those of the model: in each equation the vector of the $K$
intercepts is $\mathbf C\,(\alpha_1,\dots,\alpha_{k_2},\delta_2,\dots,\delta_{k_1})'$, and the
state-specific slopes of type $s$ are those of its state.

\subsection{The EM algorithm}\label{sec:em}

Estimation follows the EM algorithm \citep{dempster1977} as adapted to latent Markov models by
\citet{bartolucci2013}.

\paragraph{E-step.} The forward--backward recursions on the joint chain give the posterior probabilities
of the types in each period and the expected numbers of transitions between types. The posterior of the
class is the sum over the states, that of the initial state the sum over the classes, and the expected
transitions of the states are the sums of the diagonal blocks over the classes.

\paragraph{M-step.}
\begin{itemize}
\item \emph{Counts:} the weighted bivariate Poisson regression \eqref{eq:mean}--\eqref{eq:famoye} with
the posterior weights of the types, updated by Newton steps with step halving and analytic second
derivatives. The gradient and the Hessian are accumulated over the types for the expanded parameters
(common slopes, $K$ intercepts, $K$ sets of state-specific slopes, $\lambda$) and mapped to the
parameters of the model through the linear map of the intercepts. Every trial point keeps $\lambda$
inside the interval \eqref{eq:range}.
\item \emph{Dropout:} the weighted logit \eqref{eq:drop} on the periods at risk, by Newton steps with step
halving, with the same mapping.
\item \emph{Initial states and classes:} fractional multinomial logits on the posterior probabilities.
\item \emph{Transitions:} the closed form $\hat\pi_{b\mid a}\propto$ expected number of transitions from
$a$ to $b$ for the homogeneous chain; with \cmd{xtrans()} or \cmd{tvtrans}, a weighted multinomial logit
for each state of origin.
\end{itemize}
The algorithm stops when the relative change of the log likelihood falls below \cmd{emtolerance()}
(default $10^{-8}$) or after \cmd{emiterate()} iterations (default 500).

\paragraph{Direct maximization.} If the EM has not converged, or if \cmd{technique()} is specified, the
log likelihood is maximized directly with the analytic score, computed by the Fisher identity at the
cost of one forward--backward pass. Near a flat maximum the EM can stop slightly short of it (in the
simulations, by about 0.01 in log likelihood); \cmd{technique(nr)} refines the estimates.

\paragraph{Starting values.} The slopes start from Poisson regressions of the two counts and from a logit
of the dropout on the periods at risk. The state intercepts of the counts are spread around the
constant, those of the dropout are equal; the class shifts start small for the counts and spread for
the dropout; the prior logits start from zero and the chain from a persistent matrix with staying
probability 0.9. Starting values can be given with \cmd{from()}.

\paragraph{Multiple starting values.} The likelihood of a finite mixture can have local maxima, and
\citet{carrieri2026} use several starting points. With \cmd{starts(\#)} the EM is run from the default
starting values and from $\#-1$ random ones, for \cmd{startiterate()} iterations each (by default to
convergence), and goes on from the one with the largest log likelihood. The random starting values keep
the slopes, set $\lambda=0$, draw the intercepts of the states around the mean intercept, the class
shifts, the constants of the prior logits and transition matrices with staying probabilities between
0.6 and 0.95. \cmd{e(ll\_starts)} reports the log likelihood reached from each start, and
\cmd{e(n\_best)} the number of starts within 0.1 of the best.

\paragraph{Labelling.} After estimation the states are ordered by increasing intercept of the first count
equation (\cmd{sortby(eq1)}) or of the second (\cmd{sortby(eq2)}), and the classes by increasing shift
of the dropout, so that class~1 has the lowest risk of dropout. The relabelling permutes all the
parameters and does not change the likelihood.

\subsection{Variance}

The default variance is the inverse of the observed information, computed by central differences of the
analytic score \citep{mclachlan2000}. A $\lambda$ on a bound of its interval is held fixed
(\cmd{e(n\_bound)}). \cmd{vce(robust)} gives the sandwich estimator clustered on the individuals.

%=====================================================================
\section{Inference and postestimation}\label{sec:post}
%=====================================================================

\subsection{Number of states and classes}

\cmd{estat ic} reports AIC and two versions of BIC: with $\log n$, $n$ the number of individuals
(\cmd{e(bic)}, the usual choice in latent Markov models, \citealp{bartolucci2013}), and with $\log N$,
$N$ the number of person-periods (\cmd{e(bic\_N)}, the convention of Table~3 of \citealp{carrieri2026}).
It also reports the relative entropy of the classification, separately for the states and for the
classes. In applications, fit the model for a grid of $k_1$ and $k_2$, compare the BIC, and check the
size of the smallest state and class.

\subsection{Transitions and test of no transitions}

\cmd{estat transition} reports the transition matrix of the states, the expected durations
$1/(1-\pi_{a\mid a})$, the mean initial probabilities of the states and the prior probabilities of the
classes, with delta-method standard errors.

\cmd{estat lrtrans} tests $H_0$: $\pi_{b\mid a}=0$ for all $b\ne a$, the model where the states, like the
classes, are constant. $H_0$ lies on the boundary of the parameter space, and following
\citet{bartolucci2006} the likelihood-ratio statistic has a chi-bar-squared distribution, a mixture of
$\chi^2_j$, $j=0,\dots,k_2(k_2-1)$, with weights $\omega_j$ equal to the probability that the projection
of a $N(\mathbf 0,\mathbf I^{-1})$ vector onto the nonnegative orthant has $j$ positive components, where
$\mathbf I$ is the information on the off-diagonal transition probabilities net of the other parameters
at the restricted estimates. On the joint chain the derivative with respect to $\pi_{b\mid a}$ is the sum
over the classes of the derivatives with respect to the corresponding elements of the diagonal blocks.
$\mathbf I$ is estimated by the outer product of the individual scores and the weights by \cmd{reps()}
simulated projections \citep{lawson1974}. The restricted model is fitted from two sets of starting
values. The $p$-value of the $\chi^2_{k_2(k_2-1)}$ distribution, an upper bound, is also reported.

No score test is reported. The individual scores of the transition probabilities at $\bm\Pi=\mathbf I$
are sums of likelihood ratios $f_{it}(u,b)/f_{it}(u,a)$, whose distribution has very heavy tails: in one
sample, the outer product of the scores is far from the information, and a statistic
$\mathbf s'\,\widehat{\mathbf I}^{-1}\mathbf s$ is badly calibrated (Section~\ref{sec:mc}). The weights
of the chi-bar-squared depend only on the shape of $\widehat{\mathbf I}$, and the LR test keeps its
size.

\subsection{Predictions}

\cmd{predict} computes statistics that average over the latent types with one of four sets of
probabilities: posterior (given all the outcomes of the individual, the default), filtered (up to the
period), one-step-ahead (up to the period before) and prior (given the covariates and the presample
counts); \cmd{instate()} and \cmd{inclass()} fix one latent variable. The outcomes of a period are the two
counts and, in the periods at risk, the dropout. The statistics are the expected counts \cmd{mu1},
\cmd{mu2}; the probability of dropout \cmd{pdrop}; the probabilities of the states \cmd{postprob} and of
the classes \cmd{classprob}; the most likely state \cmd{state} and class \cmd{class}; the state of each
period in the most likely sequence of class and states \cmd{viterbi} \citep{viterbi1967}; the
within-type correlation \cmd{rho}; and the probability of a pair of counts \cmd{pr(}$a,b$\cmd{)}.

\subsection{Marginal effects, survival and fit}

\cmd{estat ame} reports average marginal effects on $\E(y_1)$, $\E(y_2)$ and $\Pr(D=1)$ with the
probabilities of the latent types held fixed,
\[
\mathrm{AME}=\frac1N\sum_{i,t}\sum_{s}\Pr(\text{type } s\mid\cdot)\,
\frac{m_{it}(s;A)-m_{it}(s;B)}{\Delta},
\]
with $m$ the expected count or the probability of dropout and $A$, $B$ two counterfactual designs: the
regressor shifted by $\pm h$ ($\Delta=2h$), a level and the base level of a factor variable
($\Delta=1$), one more lagged count ($\Delta=1$, through $g$) or, with \cmd{dropcounts()}, one more
count of the period in the dropout equation. State-specific slopes enter in each type. Standard errors
are by the delta method.

\cmd{estat survival} compares the Kaplan--Meier survival by period of follow-up with the survival
predicted by the model, $\prod_r(1-\bar h_r)$, with $\bar h_r$ the mean predicted probability of dropout
among the individuals at risk in period $r$ (Fig.~8 of \citealp{carrieri2026}). The default weights are
one-step-ahead, given the counts and the survival up to the period before; prior weights ignore the
selection of the survivors.

\cmd{estat fit} compares the observed distribution of each count with the distribution predicted in each
state (or class) and with their mixture (Tables~5--6 and Fig.~7 of \citealp{carrieri2026}), with the
Pearson statistic, the dissimilarity index and the largest cell difference.

\subsection{Stability of the dynamics}

\cmd{estat stability} checks whether the dynamics of the lagged counts can explode. It reports by latent
type the coefficients of the lags; with log-type lags the spectral radius of the $2\times2$ matrix of the
coefficients of $g(y_1)$ and $g(y_2)$ in the count equations, which must be below one; with linear lags
the index $\gamma_{jj}\exp(a)$ of the map $y\mapsto\exp(a+\gamma_{jj}y)$, with $a$ the 99th percentile of
the rest of the linear predictor, which must not exceed $e^{-1}$; the share of deterministic paths,
from every observation and in every type, whose expected count exceeds \cmd{limit()} within
\cmd{horizon()} periods; and the same share for simulated paths of the model from the last observed
period, with the type drawn from its posterior probabilities, together with the mean probability of
still being in the panel at the horizon.

%=====================================================================
\section{The command}\label{sec:command}
%=====================================================================

\subsection{Syntax}

\begin{lstlisting}
hmdbivpois (depvar1 = [indepvars1] [, eqopts]) (depvar2 = [indepvars2] [, eqopts])
           [(dropvar = [indepvarsd] [, eqopts])] [if] [in], id(varname) time(varname) k(#) [options]
\end{lstlisting}

The data are in long form, one observation per individual and period. \textit{indepvars},
\cmd{statevary()} and \cmd{xtrans()} may contain factor variables. Without the dropout equation the
command fits the model of \cmd{hmbivpois} with classes.

\begin{center}
\small
\begin{tabularx}{\textwidth}{lX}
\toprule
Option & Description\\
\midrule
\multicolumn{2}{l}{\emph{Equation options}}\\
\cmd{statevary(}\textit{varlist}\cmd{)} & regressors with state-specific coefficients\\
\cmd{lagvary(}\textit{depvars}\cmd{)} & lagged counts with state-specific coefficients\\
\midrule
\multicolumn{2}{l}{\emph{Model}}\\
\cmd{id()}, \cmd{time()}, \cmd{k(\#)} & individual, integer period, number of states (required)\\
\cmd{classes(\#)} & number of classes (default 1)\\
\cmd{initcov(}\textit{varlist}\cmd{)}, \cmd{classcov(}\textit{varlist}\cmd{)} & covariates of the initial states and of the classes\\
\cmd{lagtransform(}\textit{g}\cmd{)} & transformation of the lagged counts\\
\cmd{dropcounts(lag|current|both)} & counts in the dropout equation\\
\cmd{classeq(all|dropout|counts)} & equations shifted by the classes\\
\midrule
\multicolumn{2}{l}{\emph{Transitions}}\\
\cmd{xtrans(}\textit{varlist}\cmd{)}, \cmd{tvtrans}, \cmd{notransition} & covariates, free by period, none\\
\midrule
\multicolumn{2}{l}{\emph{Estimation and reporting}}\\
\cmd{starts(\#)}, \cmd{startiterate(\#)}, \cmd{startseed(\#)} & multiple starting values\\
\cmd{emiterate(\#)}, \cmd{emtolerance(\#)}, \cmd{emonly} & control of the EM\\
\cmd{technique()}, \cmd{iterate(\#)} & direct maximization\\
\cmd{from(}\textit{matname}\cmd{)} & starting values, in the order of \cmd{e(b)}\\
\cmd{sortby(eq1|eq2|none)} & labelling of the states\\
\cmd{vce(oim|robust)}, \cmd{level(\#)}, \cmd{nolog} & variance and display\\
\bottomrule
\end{tabularx}
\end{center}

Postestimation:
\begin{lstlisting}
predict [type] newvar | stub* [if] [in] [, mu1 mu2 pdrop postprob classprob state class viterbi
        rho pr(a,b) filtered onestep prior instate(#) inclass(#)]
estat ic [, n(#)]
estat transition [, period(#) level(#)]
estat ame [, prior onestep filtered instate(#) inclass(#) level(#)]
estat survival [, over(varname) onestep prior filtered posterior]
estat fit [, max(#) classes posterior onestep prior filtered]
estat lrtrans [, reps(#) seed(#)]
estat stability [, horizon(#) limit(#) nsim(#) seed(#)]
\end{lstlisting}
The stored results are listed in \cmd{help hmdbivpois}.

\subsection{Installation}

\begin{lstlisting}
net install hmdbivpois, from("https://raw.githubusercontent.com/plidonni/hmdbivpois/main/") replace
net get hmdbivpois, from("https://raw.githubusercontent.com/plidonni/hmdbivpois/main/") replace
\end{lstlisting}
The first command installs the command and its help file; the second downloads into the current folder
the example dataset \cmd{hmdbivpois\_example.dta} and the do-file \cmd{hmdbivpois\_examples.do}, which
runs all the examples below. Stata 16 or later is required.

\subsection{Examples}

The example dataset is simulated from the model with $n=2000$ individuals, a presample period and up to
$T=6$ estimation periods, $k_2=3$ states, $k_1=2$ classes and homogeneous transitions, with the true
values of Table~\ref{tab:true}. A large share of the individuals leave the panel before the end of the
follow-up, most of them in class~2, which has the higher risk of dropout. The variables \cmd{state} and \cmd{class} hold the true latent
variables.

\begin{table}[h]
\centering\footnotesize
\caption{True parameters of the example dataset}\label{tab:true}
\begin{tabular}{ll}
\toprule
Block & Values\\
\midrule
Equation $y_1$ & $x_{11}=0.4$, $x_{12}=-0.3$, $\gamma_{11}=0.03$, $\gamma_{12}=-0.05$, $\alpha_{1}=(-0.3,\,0.7,\,1.5)$, $\delta_{12}=0.3$\\
Equation $y_2$ & $x_{21}=0.3$, $x_{22}=-0.4$, $\gamma_{21}=-0.10$, $\gamma_{22}=0.10$, $\alpha_{2}=(-2.5,\,-1.5,\,-0.6)$, $\delta_{22}=0.4$\\
Dependence & $\lambda=-0.5$\\
Dropout & $x_d=0.5$, $\gamma_{d1}=0.05$, $\gamma_{d2}=0.30$, $\alpha_d=(-4.2,\,-3.7,\,-3.2)$, $\delta_{d2}=1.5$\\
Initial states (2, 3) & $z_1=(0.5,1.0)$, $y_{1,0}=(0.15,0.30)$, $y_{2,0}=(0.30,0.50)$, constant $=(-0.5,-2.0)$\\
Class 2 & $z_u=0.8$, $y_{1,0}=0.1$, $y_{2,0}=0.4$, constant $=-1.0$\\
Transitions & $\bm\Pi=\begin{pmatrix}0.85&0.12&0.03\\0.08&0.84&0.08\\0.03&0.12&0.85\end{pmatrix}$\\
\bottomrule
\end{tabular}
\end{table}

\medskip\noindent\textbf{Base model.}\par\nopagebreak
\begin{lstlisting}
use hmdbivpois_example, clear
hmdbivpois (y1 = x11 x12) (y2 = x21 x22) (died = xd), id(id) time(time) ///
    k(3) classes(2) initcov(z1) classcov(zu)
\end{lstlisting}
The output reports, for each equation, the slopes, the lag coefficients, the state intercepts
(\cmd{state1}, \dots) and the class shift (\cmd{class2}), then $\lambda$, the dropout equation, the
initial logits (\cmd{init2}, \cmd{init3}), the transition logits (\cmd{tr1\_2}, \dots) and the class logit
(\cmd{class2}), followed by the transition matrix and the shares of the states and classes.

\medskip\noindent\textbf{Postestimation.}\par\nopagebreak
\begin{lstlisting}
estat ic
estat transition
estat ame
estat survival
estat fit
predict double pd, pdrop
predict s_hat, viterbi
predict c_hat, class
tabulate state s_hat if time > 0
tabulate class c_hat if time > 0
estat lrtrans
\end{lstlisting}

\medskip\noindent\textbf{Choice of the number of states and classes.}\par\nopagebreak
\begin{lstlisting}
forvalues c = 1/2 {
    forvalues k = 2/3 {
        quietly hmdbivpois (y1 = x11 x12) (y2 = x21 x22) (died = xd), ///
            id(id) time(time) k(`k') classes(`c') initcov(z1) classcov(zu) nolog
        display "classes `c'  k `k'  BIC " %10.2f e(bic)
    }
}
\end{lstlisting}

\medskip\noindent\textbf{Dropout ignored.} Without the third equation the dropout is treated as missing at
random, and the estimates are those of \cmd{hmbivpois}; the comparison shows the bias that the
informative dropout induces.
\begin{lstlisting}
hmdbivpois (y1 = x11 x12) (y2 = x21 x22), id(id) time(time) k(3) initcov(z1)
\end{lstlisting}

\medskip\noindent\textbf{Extensions.}\par\nopagebreak
\begin{lstlisting}
hmdbivpois (y1 = x11 x12) (y2 = x21 x22) (died = xd), id(id) time(time) ///
    k(3) classes(2) initcov(z1) classcov(zu) starts(5)
matrix list e(ll_starts)

quietly hmdbivpois (y1 = x11 x12) (y2 = x21 x22) (died = xd), id(id) time(time) ///
    k(3) classes(2) initcov(z1) classcov(zu) nolog
estimates store all
quietly hmdbivpois (y1 = x11 x12) (y2 = x21 x22) (died = xd), id(id) time(time) ///
    k(3) classes(2) initcov(z1) classcov(zu) classeq(dropout) nolog
lrtest all ., df(2)

hmdbivpois (y1 = x11 x12, statevary(x11)) (y2 = x21 x22) (died = xd), id(id) time(time) ///
    k(3) classes(2) initcov(z1) classcov(zu) dropcounts(both) lagtransform(log1p)
estat stability
\end{lstlisting}
The last specification lets the slope of $x_{11}$ differ across states, puts the counts of the period in
the dropout equation and transforms the lags by $\log(1+y)$; it is misspecified for the example data and
only illustrates the syntax. Give the degrees of freedom to \cmd{lrtest} with \cmd{df()}: it takes them
from \cmd{e(rank)}, which drops when some parameters are weakly identified.

\paragraph{Practical advice.}
\begin{itemize}
\item Choose $k_2$ and $k_1$ with the BIC and check the size of the smallest state and class.
\item Use \cmd{starts()} for the final specification; when \cmd{e(n\_best)} is one, add starts. In the
simulations local maxima appeared only in overparameterized models.
\item With classes in the dropout equation only, the classes are identified by at most one dropout per
individual and the dropout parameters can be imprecise; a longer follow-up helps.
\item A $\lambda$ on a bound (\cmd{e(n\_bound)}$>0$) signals stronger dependence than the Famoye form
allows.
\item With linear lags and large counts, run \cmd{estat stability} or use a log-type
\cmd{lagtransform()}.
\end{itemize}

%=====================================================================
\section{Monte Carlo evidence}\label{sec:mc}
%=====================================================================

The command was tested with simulation experiments with $n=4000$ individuals, a presample period, $T=8$
estimation periods, $k_2=3$ states and $k_1=2$ classes, using the design of Table~\ref{tab:true} (with
about 43\% of dropout) and the changes listed in Table~\ref{tab:mc}. Each experiment has 100
replications. The table reports the ratio between the Monte Carlo standard deviation of the estimates
and the mean estimated standard error, and the coverage of the nominal 95\% confidence intervals, as
ranges over the parameters (or the marginal effects).

\begin{table}[h]
\centering\small
\caption{Monte Carlo results ($n=4000$, $T=8$, $k_2=3$, $k_1=2$, 100 replications)}\label{tab:mc}
\begin{tabularx}{\textwidth}{>{\raggedright\arraybackslash}p{5.6cm}XX}
\toprule
Design & MC s.d./mean s.e.: mean [min, max] & Coverage of 95\% CI [min, max]\\
\midrule
Base model (42 parameters) & 0.999 [0.844, 1.139] & [0.90, 1.00]\\
Average marginal effects (\cmd{estat ame}) & 1.009 [0.795, 1.158] & [0.91, 0.99]\\
Transitions on covariates (\cmd{xtrans()}, 54 parameters) & 1.007 [0.861, 1.160] & [0.89, 0.99]\\
$\log(1+y)$ lags (\cmd{lagtransform(log1p)}) & 0.997 [0.810, 1.096] & [0.91, 0.99]\\
Counts of the period in the dropout, classes in the dropout only$^a$ & 0.997 [0.642, 1.355] & [0.89, 0.99]\\
State-specific slopes in the three equations & 0.990 [0.817, 1.169] & [0.89, 0.99]\\
\bottomrule
\end{tabularx}

\smallskip
\begin{minipage}{\textwidth}\footnotesize
$^a$ Excluding the state intercepts and the class shift of the dropout equation, which have heavy-tailed
sampling distributions in this design (Monte Carlo s.d.\ 1.7--1.8, coverage 0.87--0.99): a few
replications put the probability of dropout of class~2 close to one.
\end{minipage}
\end{table}

\paragraph{Consistency checks.} Without the dropout equation and with one class, \cmd{hmdbivpois} and
\cmd{hmbivpois} give the same log likelihood, coefficients and variances (differences below $10^{-11}$),
also with transitions on covariates, every transformation of the lags and state-specific slopes. Nested
models evaluated at the restricted estimates reproduce the restricted log likelihood exactly
(\cmd{notransition}, \cmd{classeq()}, \cmd{statevary()}). With \cmd{dropcounts(both)} the estimates
equal those of the default model with the counts of the period written as regressors of the dropout.
Predictions, marginal effects in a given type and the stability diagnostics equal closed-form
computations; the Kaplan--Meier curve of \cmd{estat survival} equals the one computed directly, and the
model tracks it within about one percentage point.

\paragraph{Test of no transitions.} With data generated without transitions, the rejection rates of
\cmd{estat lrtrans} at the 10\%, 5\% and 1\% levels were 0.06, 0.05 and 0.02 with the chi-bar-squared
distribution, against 0.02, 0.01 and 0.00 with the $\chi^2_6$ bound. A score (Kuhn--Tucker) statistic
based on the outer product of the scores rejected in 0.01, 0.01 and 0.00 of the samples, and, at the true
parameters, its unprojected version had mean 44 instead of 6: it is not reported.

\paragraph{Local maxima.} In 20 samples of the base design, the best of 10 starts exceeded the default
start by at most 0.03 in log likelihood, and on average 8.6 of the 10 starts reached the best maximum. In
a model with more latent types than the data ($k_2=4$, $k_1=3$) the 10 starts ended at log likelihoods
spread over 30 points: there \cmd{starts()} matters.

%=====================================================================
\section*{Citation}
%=====================================================================
If you use \cmd{hmdbivpois}, please cite \citet{carrieri2026} and

\smallskip
\noindent Li Donni, P. 2026. hmdbivpois: Stata module to fit hidden Markov bivariate Poisson regression
with informative dropout. \url{https://github.com/plidonni/hmdbivpois}.

\begin{thebibliography}{99}
\bibitem[Bartolucci(2006)]{bartolucci2006}
Bartolucci, F. 2006. Likelihood inference for a class of latent Markov models under linear hypotheses on
the transition probabilities. \textit{Journal of the Royal Statistical Society, Series B} 68: 155--178.

\bibitem[Bartolucci and Farcomeni(2015)]{bartolucci2015}
Bartolucci, F., and A. Farcomeni. 2015. A discrete time event-history approach to informative drop-out in
mixed latent Markov models with covariates. \textit{Biometrics} 71: 80--89.

\bibitem[Bartolucci et~al.(2013)]{bartolucci2013}
Bartolucci, F., A. Farcomeni, and F. Pennoni. 2013. \textit{Latent Markov Models for Longitudinal Data}.
Boca Raton, FL: Chapman \& Hall/CRC.

\bibitem[Baum et~al.(1970)]{baum1970}
Baum, L. E., T. Petrie, G. Soules, and N. Weiss. 1970. A maximization technique occurring in the
statistical analysis of probabilistic functions of Markov chains. \textit{Annals of Mathematical
Statistics} 41: 164--171.

\bibitem[Blundell et~al.(2002)]{blundell2002}
Blundell, R., R. Griffith, and F. Windmeijer. 2002. Individual effects and dynamics in count data
models. \textit{Journal of Econometrics} 108: 113--131.

\bibitem[Carrieri et~al.(2026)]{carrieri2026}
Carrieri, V., D. Gyrd-Hansen, and P. Li Donni. 2026. Dynamic healthcare demand with informative dropout:
evidence from register data. \textit{Journal of the Royal Statistical Society Series A}.
\url{https://doi.org/10.1093/jrsssa/qnag108}.

\bibitem[Dempster et~al.(1977)]{dempster1977}
Dempster, A. P., N. M. Laird, and D. B. Rubin. 1977. Maximum likelihood from incomplete data via the EM
algorithm. \textit{Journal of the Royal Statistical Society, Series B} 39: 1--38.

\bibitem[Famoye(2010)]{famoye2010}
Famoye, F. 2010. On the bivariate negative binomial regression model. \textit{Journal of Applied
Statistics} 37: 969--981.

\bibitem[Heckman(1981)]{heckman1981}
Heckman, J. J. 1981. The incidental parameters problem and the problem of initial conditions in
estimating a discrete time--discrete data stochastic process. In \textit{Structural Analysis of
Discrete Data with Econometric Applications}, ed. C. F. Manski and D. McFadden, 179--195. Cambridge, MA:
MIT Press.

\bibitem[Lakshminarayana et~al.(1999)]{lakshminarayana1999}
Lakshminarayana, J., S. N. N. Pandit, and K. Srinivasa Rao. 1999. On a bivariate Poisson distribution.
\textit{Communications in Statistics -- Theory and Methods} 28: 267--276.

\bibitem[Laudicella and Li Donni(2022)]{laudicella2022}
Laudicella, M., and P. Li Donni. 2022. The dynamic interdependence in the demand of primary and emergency
secondary care: A hidden Markov approach. \textit{Journal of Applied Econometrics} 37: 521--536.

\bibitem[Lawson and Hanson(1974)]{lawson1974}
Lawson, C. L., and R. J. Hanson. 1974. \textit{Solving Least Squares Problems}. Englewood Cliffs, NJ:
Prentice-Hall.

\bibitem[McLachlan and Peel(2000)]{mclachlan2000}
McLachlan, G., and D. Peel. 2000. \textit{Finite Mixture Models}. New York: Wiley.

\bibitem[Viterbi(1967)]{viterbi1967}
Viterbi, A. J. 1967. Error bounds for convolutional codes and an asymptotically optimum decoding
algorithm. \textit{IEEE Transactions on Information Theory} 13: 260--269.

\bibitem[Wooldridge(2005)]{wooldridge2005}
Wooldridge, J. M. 2005. Simple solutions to the initial conditions problem in dynamic, nonlinear panel
data models with unobserved heterogeneity. \textit{Journal of Applied Econometrics} 20: 39--54.

\bibitem[Zucchini et~al.(2016)]{zucchini2016}
Zucchini, W., I. L. MacDonald, and R. Langrock. 2016. \textit{Hidden Markov Models for Time Series: An
Introduction Using R}. 2nd ed. Boca Raton, FL: Chapman \& Hall/CRC.
\end{thebibliography}

\end{document}
